\chapter{Dataframes}\label{ch:pl:dataframes}

Deux chercheurs calculent des barres de cinq minutes à partir des mêmes transactions et obtiennent des volumes différents. Une bibliothèque nomme une barre par son left
edge et place une transaction horodatée exactement sur une borne dans la barre qui commence à cette borne~; l'autre la nomme par son right edge et
place la même transaction dans la barre qui se termine à cette borne~; l'un des deux supprime les minutes pendant lesquelles rien n'a été échangé, l'autre non.
Les deux sont justes selon leurs propres conventions, et un signal qui compare les deux ensembles de barres trouve une régularité qui n'est rien d'autre que la
convention. Le \omterm{def:pl:dataframes:dataframe}{dataframe} est le langage de travail de la plateforme de recherche~; ses idiomes --- barres, jointures as-of, fenêtres glissantes,
rangs cross-sectionnels, pivots --- c'est là que se cachent ses bugs silencieux, et ses moteurs diffèrent énormément par leur coût. Ce chapitre
formule une fois pour toutes ces idiomes avec des conventions explicites pour trois moteurs (\texttt{firm.marketdf}), vérifie que les moteurs concordent, et
exécute sur un mois de cotations et de transactions un même pipeline de recherche, en \omterm{def:pl:dataframes:eager}{évaluation immédiate}, en \omterm{def:pl:dataframes:eager}{évaluation paresseuse} et en \omterm{def:pl:dataframes:outofcore}{traitement hors mémoire}.

\section{Le modèle du dataframe}

\begin{definition}[Dataframe]\label{def:pl:dataframes:dataframe}
Un \emph{dataframe}\index{dataframe} est une table de colonnes nommées et typées, de même longueur, munie d'opérations qui agissent sur des colonnes entières
et sur des \omterm{def:pl:columnar-formats:rowgroup}{groupes de lignes}~: sélection, filtrage, colonnes calculées, regroupement et agrégation, jointures, tri, fenêtres et remodelage.
\end{definition}

Les trois moteurs du chapitre mettent en œuvre le même modèle de façon différente. pandas conserve chaque colonne sous forme de tableau numpy (ou
tableau Arrow) muni d'un index de lignes, et exécute chaque opération dès qu'elle est appelée. Polars conserve les colonnes dans la disposition mémoire d'Arrow,
n'a pas d'index de lignes, et peut soit exécuter chaque opération, soit construire un plan et l'exécuter plus tard. DuckDB est un moteur SQL~: une requête
est un plan, exécuté lorsque son résultat est demandé (chapitre~5). La \omterm{def:pl:columnar-formats:layouts}{disposition en colonnes} du chapitre~3 sous-tend les trois.

\begin{definition}[Format long, format large]\label{def:pl:dataframes:shape}
Une table en \emph{format long}\index{format long} a une ligne par observation et une colonne pour chaque variable, y compris les clés
(symbole, temps, valeur). Les mêmes données en \emph{format large}\index{format large} ont une ligne par valeur d'une clé et une colonne par
valeur de l'autre (une ligne par minute, une colonne par symbole).
\end{definition}

\begin{omfigure}
\begin{tikzpicture}[font=\scriptsize,cell/.style={draw,minimum width=1.05cm,minimum height=0.42cm,inner sep=1pt}]
\foreach \r/\t/\s/\v in {0/09:31/AAA/10.1,1/09:31/BBB/20.4,2/09:32/AAA/10.2,3/09:32/BBB/20.3} {
  \node[cell] at (0,-\r*0.42) {\t}; \node[cell] at (1.05,-\r*0.42) {\s}; \node[cell] at (2.1,-\r*0.42) {\v};}
\foreach \x/\h in {0/heure,1.05/symbole,2.1/clôture} {\node[cell,fill=omDef!15,font=\tiny] at (\x,0.42) {\h};}
\node[font=\scriptsize\itshape] at (1.05,1.05) {long~: une ligne par observation};
\draw[->,thick] (3.1,-0.4) -- node[above]{pivot} (4.6,-0.4);
\draw[<-,thick] (3.1,-0.9) -- node[below]{dépivoter} (4.6,-0.9);
\foreach \r/\t/\a/\b in {0/09:31/10.1/20.4,1/09:32/10.2/20.3} {
  \node[cell] at (5.6,-\r*0.42) {\t}; \node[cell] at (6.65,-\r*0.42) {\a}; \node[cell] at (7.7,-\r*0.42) {\b};}
\foreach \x/\h in {5.6/heure,6.65/AAA,7.7/BBB} {\node[cell,fill=omMeth!15] at (\x,0.42) {\h};}
\node[font=\scriptsize\itshape,align=center] at (6.65,1.2) {large~: une ligne par minute,\\une colonne par symbole};
\end{tikzpicture}

\omcaption{Les mêmes cours de clôture en \omterm{def:pl:dataframes:shape}{format long} (clés en colonnes) et en \omterm{def:pl:dataframes:shape}{format large} (une clé en lignes, l'autre en colonnes).
Le code de recherche passe de l'un à l'autre~: le \omterm{def:pl:dataframes:shape}{format long} pour le stockage, les jointures et les regroupements~; le \omterm{def:pl:dataframes:shape}{format large} pour l'arithmétique cross-sectionnelle et
les matrices.}\label{fig:pl:dataframes:shape}
\end{omfigure}

Les données de marché sont stockées au \omterm{def:pl:dataframes:shape}{format long} (les partitions du chapitre~4 sont des tables longues), et la plupart des transformations sont plus claires sous cette forme~: une fenêtre glissante
est une fenêtre sur les lignes d'un même symbole, un rang cross-sectionnel est un rang sur les lignes d'une même minute. Le \omterm{def:pl:dataframes:shape}{format large} sert aux
opérations matricielles de la construction de portefeuille (livre~7, chapitres~24--26), et le pivot entre les deux est l'endroit où
apparaissent les valeurs manquantes~: une minute où un symbole n'a pas été échangé est une cellule vide du tableau large, jamais une ligne du tableau long.

\section{Évaluation immédiate et évaluation paresseuse}

\begin{definition}[Évaluation immédiate, évaluation paresseuse]\label{def:pl:dataframes:eager}
Un moteur à \emph{évaluation immédiate}\index{évaluation immédiate} (eager evaluation) exécute chaque opération dès qu'elle est appelée et matérialise l'intégralité
de son résultat. Un moteur à \emph{évaluation paresseuse}\index{évaluation paresseuse} (lazy evaluation) enregistre les opérations sous forme de plan et n'exécute ce plan que
lorsqu'un résultat est demandé, ce qui lui permet de voir l'ensemble du calcul avant d'en exécuter la moindre partie.
\end{definition}

L'\omterm{def:pl:dataframes:eager}{évaluation immédiate} est facile à raisonner et à déboguer pas à pas~; chaque table intermédiaire existe, et chaque table intermédiaire
coûte de la mémoire. L'\omterm{def:pl:dataframes:eager}{évaluation paresseuse} donne au moteur l'ensemble du pipeline, ce qui rend l'optimisation possible.

\section{Plans de requête et optimisation}

\begin{definition}[Optimiseur de requêtes]\label{def:pl:dataframes:optimiser}
Un \emph{optimiseur de requêtes}\index{optimiseur de requêtes} (query optimiser) réécrit un plan en un plan équivalent et moins coûteux avant de l'exécuter~: il reporte les filtres
et les projections vers les lectures (pour que moins de lignes et de colonnes soient lues, chapitre~3), supprime les colonnes et les opérations inutiles, réordonne les
jointures et les agrégations, et choisit les algorithmes.
\end{definition}

Une requête en \omterm{def:pl:dataframes:eager}{évaluation paresseuse} sur les cours d'un symbole pour un jour donné, écrite comme un parcours de chaque fichier de transactions suivi d'un filtre et d'une sélection, est
optimisée par Polars en un plan qui ne lit qu'une partition et quatre des six colonnes, et applique le filtre pendant le parcours
(la \cref{fig:pl:dataframes:plan}).

\begin{omfigure}
\begin{tikzpicture}[font=\scriptsize,op/.style={draw,rounded corners=2pt,align=center,minimum width=4.4cm,minimum height=0.62cm,
  inner sep=2pt},>=Stealth]
\node[font=\small\bfseries] at (0,0.75) {tel qu'écrit};
\node[op,fill=omMeth!12] (a1) at (0,0) {sélection \texttt{ts}, \texttt{price}};
\node[op,fill=omMeth!12] (a2) at (0,-1.0) {filtre symbole $=$ S007, date $=5$};
\node[op,fill=omThm!12] (a3) at (0,-2.0) {parcours des 20 partitions par date\\ les 6 colonnes};
\draw[->] (a3) -- (a2); \draw[->] (a2) -- (a1);
\draw[->,very thick,gray] (2.6,-1.0) -- node[above,black]{optimiseur} (4.2,-1.0);
\node[font=\small\bfseries] at (6.8,0.75) {tel qu'exécuté};
\node[op,fill=omMeth!12] (b1) at (6.8,0) {sélection \texttt{ts}, \texttt{price} (2 sur 2)};
\node[op,fill=omDef!12,minimum height=1.62cm] (b2) at (6.8,-1.5) {parcours de \texttt{date=5} seulement\\ projection de 4 colonnes sur 6\\ filtre symbole $=$ S007\\
  pendant le parcours};
\draw[->] (b2) -- (b1);
\node[anchor=north,align=center,text width=4.6cm] at (6.8,-2.45) {élagage des partitions, report des projections et report des prédicats};
\end{tikzpicture}

\omcaption{Une requête restreinte telle qu'écrite et telle que Polars l'exécute (d'après le plan qu'il affiche). Le filtre sur la colonne de partition supprime 19 partitions sur
20 avant qu'aucun fichier ne soit ouvert~; la projection conserve les deux colonnes renvoyées et les deux que lit le filtre~; le filtre sur le
symbole s'exécute à l'intérieur du parcours.}\label{fig:pl:dataframes:plan}
\end{omfigure}

L'élagage des partitions (seule \texttt{date=5} est parcourue), le \omterm{def:pl:columnar-formats:pushdown}{report des projections} (quatre colonnes~: les deux demandées et les deux dont le filtre
a besoin) et le \omterm{def:pl:columnar-formats:pushdown}{report des prédicats} du chapitre~3 sont tous du ressort de l'optimiseur~; une version en \omterm{def:pl:dataframes:eager}{évaluation immédiate} du même code lirait le mois entier et
le filtrerait en mémoire.

\section{Traitement hors mémoire et exécution en flux}

\begin{definition}[Traitement hors mémoire, exécution en flux]\label{def:pl:dataframes:outofcore}
Le \emph{traitement hors mémoire}\index{traitement hors mémoire} (out-of-core processing) calcule un résultat sur des données plus grandes que la mémoire en les lisant par morceaux
depuis le stockage. L'\emph{exécution en flux}\index{exécution en flux} (streaming execution) en est la forme au niveau du moteur~: les opérateurs du plan se transmettent des lots
au fur et à mesure de leur production, de sorte qu'un parcours, un filtre et une agrégation ne retiennent qu'un lot à la fois, tandis que les opérateurs qui ont besoin
de la totalité de leur entrée (un tri global, certaines jointures) la mettent en tampon.
\end{definition}

Le chapitre~8 écrivait à la main des boucles de \omterm{def:pl:dataframes:outofcore}{traitement hors mémoire}~; un moteur à \omterm{def:pl:dataframes:eager}{évaluation paresseuse} (lazy) doté d'un exécuteur en flux fait de même pour un plan entier. Il ne peut pas le faire
pour tous les opérateurs~: une jointure as-of exige que ses deux côtés soient triés, et un tri a besoin de toute son entrée, si bien qu'un plan qui en contient met en tampon à ces
points. Le pipeline ci-dessous en contient les deux, ce qui explique que l'\omterm{def:pl:dataframes:outofcore}{exécution en flux} économise ici de la mémoire sans pour autant la rendre constante.

\section{Idiomes pour les données de marché}

\begin{definition}[Fonction de fenêtre]\label{def:pl:dataframes:window}
Une \emph{fonction de fenêtre}\index{fonction de fenêtre} calcule, pour chaque ligne, une valeur sur un ensemble de lignes apparentées --- les lignes d'un même
groupe, ordonnées, éventuellement limitées à un cadre comme les $n$ dernières lignes --- sans réduire les lignes comme le fait une agrégation~:
une moyenne glissante par symbole, un rang au sein d'une minute, la valeur de la ligne précédente (un décalage).
\end{definition}

Cinq idiomes couvrent l'essentiel de la recherche sur les données de marché, et chacun repose sur une convention qu'il faut énoncer explicitement plutôt que de la reprendre telle quelle.

\begin{method}[Idiomes des données de marché et leurs conventions]\label{met:pl:dataframes:idioms}
\textbf{Barres}~: indiquer la largeur, l'étiquette (bord gauche ou droit), le côté fermé, et si les barres vides sont supprimées ou comblées (avec
le cours de clôture précédent et un volume nul). \textbf{Jointures as-of}~: indiquer la clé (le temps ou le numéro de séquence, chapitre~4), le groupe et si
une égalité de temps correspond. \textbf{Fenêtres glissantes}~: indiquer le cadre (lignes ou temps), le groupe et le nombre minimal d'observations.
\textbf{Rangs en coupe transversale}~: indiquer la règle de départage. \textbf{Tableaux croisés}~: indiquer ce que signifie une cellule manquante.
\end{method}

\omcode{code/firm/marketdf/firm_marketdf.py}{54}{88}{Barres avec conventions explicites dans trois moteurs~: l'indice de barre, déduit du
côté fermé par division plancher (quel que soit le comportement de la division entière du moteur sur les nombres négatifs), l'étiquette déduite du bord, puis la même
agrégation.}\label{lst:pl:dataframes:bars}

Le \Cref{lst:pl:dataframes:bars} illustre l'idiome des barres dans les trois moteurs. Deux détails sont de ceux qui rompent l'équivalence en silence~: le
premier et dernier cours d'une barre doivent être pris dans l'ordre temporel avec un départageur (le numéro de séquence du flux), et l'indice d'une
barre fermée à droite est $\lceil t/w \rceil - 1 = \lfloor (t-1)/w \rfloor$ pour des temps entiers, et chaque moteur doit le calculer par division plancher
--- la division entière de DuckDB tronque vers zéro, et la version SQL calcule le plancher explicitement. Les tests d'acceptation
comparent les trois moteurs sur des transactions horodatées à la seconde, de sorte que les bornes et les égalités apparaissent.

\begin{example}[Deux jeux de barres]\label{ex:pl:dataframes:bars}
Sur la journée d'un même symbole de transactions synthétiques (375 transactions, horodatées à la seconde comme certains fournisseurs les livrent), des barres de cinq minutes étiquetées
et fermées à gauche ou à droite diffèrent en volume dans 4 des 74 intervalles~: deux transactions tombent exactement sur une borne de cinq minutes,
et chacune déplace son volume, soit 600 actions au total, d'une barre à la suivante. Les deux jeux portent aussi des étiquettes distantes de cinq minutes pour le même
intervalle. Combler les barres vides ajoute 4 barres aux 74, et les rendements par barre de la journée passent à 77 au lieu de 73~: un rendement calculé à cheval sur une
barre supprimée couvre dix minutes et se compare, dans une fenêtre glissante, à des rendements de cinq minutes.
\end{example}

\begin{omfigure}
\begin{tikzpicture}[font=\scriptsize,>=Stealth]
\draw[->] (0,0) -- (11.2,0) node[right]{temps};
\foreach \x/\l in {0.5/10:00,5.5/10:05,10.5/10:10} {\draw (\x,0.1) -- (\x,-0.1) node[below]{\l};}
\fill[omThm] (5.5,0) circle (2.5pt);
\node[omThm,above] at (5.5,0.1) {transaction à 10:05:00};
\draw[omDef,thick] (0.5,0.9) -- (5.4,0.9); \draw[omDef,line width=2pt] (5.6,0.9) -- (10.5,0.9);
\fill[omDef] (0.5,0.9) circle (2pt); \draw[omDef,fill=white] (5.4,0.9) circle (2pt);
\fill[omDef] (5.6,0.9) circle (2pt); \draw[omDef,fill=white] (10.5,0.9) circle (2pt);
\node[omDef,above] at (3,0.95) {barre étiquetée 10:00}; \node[omDef,above,align=center] at (8,0.95) {barre étiquetée 10:05\\(qui contient la transaction)};
\node[omDef,anchor=east,font=\scriptsize\bfseries] at (0.3,0.9) {gauche/gauche};
\draw[omMeth,line width=2pt] (0.5,-0.9) -- (5.4,-0.9); \draw[omMeth,thick] (5.6,-0.9) -- (10.5,-0.9);
\draw[omMeth,fill=white] (0.5,-0.9) circle (2pt); \fill[omMeth] (5.4,-0.9) circle (2pt);
\draw[omMeth,fill=white] (5.6,-0.9) circle (2pt); \fill[omMeth] (10.5,-0.9) circle (2pt);
\node[omMeth,below,align=center] at (3,-0.95) {barre étiquetée 10:05\\(qui contient la transaction)}; \node[omMeth,below] at (8,-0.95) {barre étiquetée 10:10};
\node[omMeth,anchor=east,font=\scriptsize\bfseries] at (0.3,-0.9) {droite/droite};
\end{tikzpicture}

\omcaption{Une transaction horodatée exactement à 10:05:00 selon deux conventions de barres. Fermée à gauche, elle appartient à la barre $[10\!:\!05,
10\!:\!10)$, étiquetée 10:05~; fermée à droite, elle appartient à la barre $(10\!:\!00, 10\!:\!05]$, également étiquetée 10:05 parce que cette convention nomme les barres par
leur bord droit. Les points pleins marquent les extrémités fermées, les points vides les extrémités ouvertes~; la barre épaisse contient la transaction. La même étiquette désigne deux
intervalles différents, et le volume de la transaction change d'intervalle.}\label{fig:pl:dataframes:edge}
\end{omfigure}

\section{Jointures, fenêtres et rangs~: là où se cachent les égalités}

Les barres sont la convention la plus visible~; trois autres se cachent dans les jointures et les fenêtres. Une jointure as-of sur l'horodatage associe chaque transaction
à la dernière cotation à cet instant ou avant~: lorsqu'une transaction et la cotation qu'elle a modifiée partagent un horodatage, la jointure inclusive renvoie
la cotation postérieure à la transaction. Le chapitre~4 l'a relevé sur 223 des 17\,281 transactions et a rendu la jointure stricte sur le numéro de séquence du flux~;
\texttt{firm\_marketdf.asof} prend \texttt{strict} comme argument dans les trois moteurs (pandas \texttt{allow\_exact\_matches}, Polars
de même, DuckDB l'inégalité \texttt{>} au lieu de \texttt{>=}).

Une fenêtre glissante a un cadre en lignes ou en temps, et un nombre minimal d'observations en deçà duquel elle ne renvoie rien ou une
valeur partielle~: la moyenne glissante de cet idiome porte sur toutes les lignes existantes, soit une seule au début de chaque symbole, ce qui doit être connu du code qui
l'utilise. Un rang en coupe transversale a une règle de départage, et les moteurs diffèrent par défaut.

\begin{example}[Rendements à égalité]\label{ex:pl:dataframes:ties}
Sur une minute, quatre symboles ont des rendements de $0.1$, $0.2$, $0.2$ et $0.3$ (en points de base). La règle de la moyenne donne les rangs 1,
2.5, 2.5 et 4, la somme des rangs $10 = 4\cdot5/2$ quelles que soient les égalités. Le \texttt{rank()} de SQL donne 1, 2, 2, 4 (la règle du minimum), ce qui abaisse
le rang moyen des symboles à égalité~; avec $c$ lignes à égalité pour \texttt{rank()} $= r$, le rang moyen vaut $(2r + c - 1)/2 = (4 + 2 - 1)/2 = 2.5$,
l'expression qu'utilise le SQL du \cref{lst:pl:dataframes:sql}. Les rendements de cours sur une grille de pas de cotation présentent souvent des égalités~: dans le mois le plus court des
tests, 10.8~\% des 222\,621 rendements à une minute sont exactement nuls, et 6\,467 groupes de rendements à égalité apparaissent dans 7\,808 minutes.
\end{example}

\section{Un seul pipeline, cinq modes}

Le pipeline de recherche du chapitre s'exécute sur le mois de cotations du chapitre~3 (désormais 3 millions, soit trois mille par symbole et par jour) et
un mois de 1.2 million de transactions, stocké en partitions par date au format Parquet~: des barres à une minute, leurs rendements logarithmiques, le rang en coupe transversale de
chaque rendement au sein de sa minute, le rang moyen de chaque symbole sur le mois~; et, en parallèle, chaque transaction rattachée par jointure as-of à sa
cotation en vigueur, et la fourchette effective moyenne de chaque symbole, en points de base. Le \Cref{lst:pl:dataframes:lazy} est la version Polars, un plan en
\omterm{def:pl:dataframes:eager}{évaluation paresseuse} que partagent les trois modes Polars.

\omcode{code/platforms/10-dataframes/python/pl_dataframes.py}{111}{124}{Le pipeline sous forme d'un seul plan Polars~: barres, rendements et rangs par
expressions de fenêtre, la jointure as-of et la fourchette effective~; le même plan est collecté en évaluation immédiate (eager), en évaluation paresseuse ou par le
moteur en flux.}\label{lst:pl:dataframes:lazy}

\omcode{code/platforms/10-dataframes/python/pl_dataframes.py}{140}{161}{Le même pipeline en SQL de DuckDB, sur les fichiers Parquet~:
barres par agrégat ordonné, rendements par décalage, rangs par une fenêtre dont la règle de départage est explicitée, et la
jointure as-of.}\label{lst:pl:dataframes:sql}

Les cinq modes renvoient les mêmes cinquante rangs moyens et cinquante fourchettes moyennes, à neuf décimales près. La \Cref{fig:pl:dataframes:engines} les mesure
chacun dans son propre processus, avec un seul thread, la mémoire étant comptée comme le high-water mark du processus moins celui d'un processus qui ne fait que
importer les bibliothèques (155~Mo). pandas, en \omterm{def:pl:dataframes:eager}{évaluation immédiate} et avec un index de lignes, lit tout le mois en mémoire et conserve chaque table intermédiaire~:
il a besoin de 895~Mo de plus que l'import des bibliothèques et de 1.3~s. Polars en \omterm{def:pl:dataframes:eager}{évaluation immédiate} conserve les mêmes tables sous forme Arrow~: 769~Mo. En \omterm{def:pl:dataframes:eager}{évaluation paresseuse}, il lui faut 654~Mo, et
avec le moteur en flux 614~Mo~; la jointure as-of et les tris limitent le gain à un niveau modeste. DuckDB, qui exécute le SQL sur les fichiers avec ses
propres opérateurs en flux, a besoin de 317~Mo, soit un tiers de pandas.

\begin{omfigure}
\begin{tikzpicture}
\begin{axis}[width=12cm,height=5.6cm,ybar,bar width=10pt,ymin=0,ymax=1000,ylabel={mémoire supplémentaire (Mo)},
  xtick={0,1,2,3,4},xticklabels from table={figdata/platforms/10-dataframes/measured_engines.csv}{engine},
  tick label style={font=\small},xticklabel style={text width=1.7cm,align=center},label style={font=\small},grid=major,grid style={gray!20},enlarge x limits=0.15,
  nodes near coords,point meta=explicit symbolic,nodes near coords style={font=\footnotesize},table/col sep=comma]
\addplot[fill=omDef!45,draw=omDef] table[x expr=\coordindex,y=extra_mb,meta=seconds]{figdata/platforms/10-dataframes/measured_engines.csv};
\end{axis}
\end{tikzpicture}

\omcaption{Mémoire au-delà des imports (barres) et temps réel en secondes (étiquettes) du même pipeline de recherche sur un mois de cotations et
de transactions, dans cinq modes de moteur, chacun exécuté dans son propre processus avec un seul fil. Les cinq donnent des réponses identiques. Mesuré sur un ordinateur portable (Intel
Core Ultra 7 155H) sous WSL2, machine par ailleurs inactive. Données~: \texttt{bench\_dataframes.py}.}\label{fig:pl:dataframes:engines}
\end{omfigure}

\begin{remark}[Mesurer la mémoire d'un processus fils]\label{rem:pl:dataframes:rss}
La première version de la mesure lisait le pic de l'ensemble résident du fils dans \texttt{getrusage}, et chaque moteur semblait demander à peu près
la même mémoire~: sous Linux, les statistiques d'utilisation sont conservées à travers \texttt{exec}, si bien que le chiffre peut inclure la mémoire du processus à partir duquel le fils a été
lancé. La high-water mark de
\texttt{/proc/self/status} appartient à l'espace d'adressage propre du processus~; c'est elle que le chapitre rapporte.
\end{remark}

\begin{dated}{2026-09}{Versions des moteurs}\label{dat:pl:dataframes:versions}
Les mesures et les plans du chapitre utilisent pandas 2.3.3, Polars 1.44.2, DuckDB 1.5.5 et PyArrow 25.0.1, les versions figées dans les
prérequis de la série. Les plans, les valeurs par défaut et l'usage mémoire changent d'une version à l'autre~: le plan de la \cref{fig:pl:dataframes:plan} est celui qu'affiche cette
version de Polars, et Polars se replie sur son moteur en mémoire pour les opérations pas encore implémentées en flux.
\end{dated}

\section{Tutoriel~: barres, conventions et cinq moteurs}\label{tut:pl:dataframes:tut}

\begin{tutorial}
\textbf{Objectif.} Calculer des barres selon des conventions énoncées dans trois moteurs, observer ce que changent les conventions, et exécuter un pipeline dans cinq
modes avec des réponses identiques. \textbf{État final~:} l'\cref{ex:pl:dataframes:bars}, la \cref{fig:pl:dataframes:engines} et le plan optimisé
du chapitre.
\begin{tutsteps}
\item \textbf{Idiomes}~: \texttt{firm\_marketdf.bars}, \texttt{asof}, \texttt{rolling\_mean}, \texttt{xs\_rank}, \texttt{to\_wide},
\texttt{to\_long}~; les tests d'acceptation vérifient que les trois moteurs concordent.
\item \textbf{Conventions}~: \texttt{pl\_dataframes.conventions()} pour la journée d'un symbole de l'\cref{ex:pl:dataframes:bars}.
\item \textbf{Plan}~: \texttt{plan\_text()} affiche le plan paresseux optimisé d'une requête étroite.
\item \textbf{Pipeline}~: \texttt{write\_month(**BIG)}, puis \texttt{bench\_dataframes.py} exécute les cinq modes, chacun dans son propre processus,
vérifie les réponses et enregistre le temps et la mémoire.
\end{tutsteps}
\textbf{Pour aller plus loin.} Remplacer la jointure as-of du plan Polars par une jointure sur la seconde précédente et observer ce que le moteur en flux
économise alors~; ajouter une option de barres remplies au pipeline et mesurer comment les rangs moyens changent.
\end{tutorial}

\section{Construction~: idiomes des données de marché}\label{bld:pl:dataframes:marketdf}

\begin{build}
\bfield{Rôle} L'implémentation partagée par la plateforme des idiomes que chaque chercheur réécrit~: un seul contrat, trois moteurs, les
conventions en arguments, l'équivalence en tests. Le moteur de backtest du chapitre~11 et le code de variables de l'environnement de recherche
(chapitre~15) l'utilisent.
\bfield{Interface} \texttt{ENGINES}~; \texttt{bars(ticks, width, label, closed, fill, engine)}~;
\texttt{asof(left, right, by, on, engine, strict)}~; \texttt{rolling\_mean(df, by, order, column, window, engine)}~; \texttt{xs\_rank(df, time, column, engine)}~;
\texttt{to\_wide}, \texttt{to\_long}.
\bfield{Règles} Les conventions sont des arguments sans valeur par défaut cachée au-delà de celle qui est documentée~; chaque idiome donne le même résultat dans
chaque moteur, ligne pour ligne~; les égalités sont départagées par le numéro de séquence~; les temps entiers sont regroupés par division entière arrondie vers le bas dans chaque moteur.
\bfield{Tests d'acceptation} \texttt{code/firm/marketdf/tests/}~: barres selon quatre conventions égales entre moteurs sur des transactions datées à la seconde,
avec volume conservé~; les conventions vérifiées à la main sur quatre transactions~; les jointures as-of (inclusives et strictes), moyennes mobiles et rangs égaux
entre moteurs~; un aller-retour de pivot.
\bfield{Pour aller plus loin} Fenêtres glissantes en temps~; barres à prix pondérés par le volume~; les idiomes pour les tables Arrow sans pandas.
\end{build}

\begin{omsources}
\item H.~Wickham, ``Tidy data'', \emph{Journal of Statistical Software} 59(10), 2014.
\item Documentation de pandas, \texttt{DataFrame.resample}~; guide utilisateur de Polars, ``Streaming''~; documentation de DuckDB, ``AsOf join''.
\item Pages de manuel Linux, \texttt{getrusage(2)}.
\end{omsources}

\section{Exercices}

\begin{exercise}[$\star$]\label{exo:pl:dataframes:1}
Une transaction est horodatée à 10:05:00 exactement. Dans quelle barre de cinq minutes, et sous quelle étiquette, tombe-t-elle selon chacune des deux conventions
de la \cref{fig:pl:dataframes:edge}~?
\end{exercise}

\begin{exercise}[$\star$]\label{exo:pl:dataframes:2}
Convertir en \omterm{def:pl:dataframes:shape}{format large} le tableau au \omterm{def:pl:dataframes:shape}{format long} de la \cref{fig:pl:dataframes:shape}, puis revenir en arrière. Que contient le \omterm{def:pl:dataframes:shape}{format large} pour une minute où BBB
n'a pas été échangé~?
\end{exercise}

\begin{exercise}[$\star$]\label{exo:pl:dataframes:3}
Quelles colonnes et quelles partitions le plan optimisé de la requête étroite lit-il, et pourquoi quatre colonnes plutôt que deux~?
\end{exercise}

\begin{exercise}[$\star\star$]\label{exo:pl:dataframes:4}
Pourquoi l'indice d'une barre fermée à droite est-il $\lfloor (t-1)/w \rfloor$ pour un temps entier $t$, et que se passe-t-il dans un moteur dont la division entière
tronque vers zéro, pour des temps négatifs~?
\end{exercise}

\begin{exercise}[$\star\star$]\label{exo:pl:dataframes:5}
Dans l'\cref{ex:pl:dataframes:bars}, supprimer les barres vides donne 73 rendements et les remplir en donne 77. Quels rendements diffèrent, et que surestime une
volatilité glissante sur 20 barres de la série où les barres vides ont été supprimées~?
\end{exercise}

\begin{exercise}[$\star\star$]\label{exo:pl:dataframes:6}
Pourquoi l'\omterm{def:pl:dataframes:outofcore}{exécution en flux} économise-t-elle moins de mémoire sur le pipeline du chapitre que sur un pipeline de parcours, de filtres et d'agrégations~?
\end{exercise}

\begin{exercise}[$\star\star\star$]\label{exo:pl:dataframes:7}
\emph{Programmation.} Avec \texttt{firm\_marketdf.bars}, calculer les barres d'une minute du mois du chapitre pour un symbole dans les trois moteurs et
selon les quatre conventions des tests d'acceptation. Les moteurs concordent-ils, et le volume total dépend-il de la convention~?
\end{exercise}

\begin{exercise}[$\star\star\star$]\label{exo:pl:dataframes:8}
\emph{Trouver la faille.} «~Nous rééchantillonnons les transactions en barres avec la valeur par défaut de la bibliothèque, calculons les rendements avec \texttt{pct\_change} sur la table de barres
triée par temps, puis classons-les entre symboles à chaque minute.~»
\end{exercise}

\section{Problème~: deux jeux de barres}

\begin{problem}[{Problème du week-end --- un mois à travers cinq moteurs}]\label{pb:pl:dataframes:1}
Le mois de cotations et de transactions du chapitre, ses conventions de barres et les cinq modes de moteur de la \cref{fig:pl:dataframes:engines}.

\medskip\noindent\textbf{Partie I --- Conventions.}
\begin{enumerate}
\item Combien de transactions de la journée d'un symbole tombent exactement sur une frontière de cinq minutes, et pourquoi cela peut-il arriver~?
\item Combien des 74 intervalles diffèrent en volume entre les deux conventions, et combien d'actions se déplacent~?
\item À quelle distance sont les étiquettes d'un même intervalle selon les deux conventions~?
\item Combien de barres et de rendements compte la journée lorsque les barres vides sont supprimées, et combien lorsqu'elles sont remplies~?
\item Quelle convention retenir pour un entrepôt de recherche, et que faut-il consigner~?
\end{enumerate}

\medskip\noindent\textbf{Partie II --- Plans.}
\begin{enumerate}[resume]
\item Que matérialise un moteur à \omterm{def:pl:dataframes:eager}{évaluation immédiate} qu'un moteur à \omterm{def:pl:dataframes:eager}{évaluation paresseuse} n'a pas besoin de matérialiser~?
\item Quelles sont les trois optimisations que révèle le plan de la requête restreinte~?
\item Pourquoi un plan à \omterm{def:pl:dataframes:eager}{évaluation paresseuse} est-il plus facile à optimiser qu'un code à \omterm{def:pl:dataframes:eager}{évaluation immédiate}~?
\item Quels opérateurs du pipeline ne peuvent pas s'exécuter en flux, et pourquoi~?
\item Pourquoi le SQL doit-il calculer explicitement la division à l'arrondi inférieur~?
\end{enumerate}

\medskip\noindent\textbf{Partie III --- Les cinq modes.}
\begin{enumerate}[resume]
\item Que calcule le pipeline, et comment compare-t-on les cinq résultats~?
\item De combien de mémoire, en plus des imports, chaque mode a-t-il besoin~?
\item Quel mode est le plus rapide, et lequel utilise le moins de mémoire~?
\item Pourquoi pandas est-il le plus coûteux en mémoire~?
\item Pourquoi une première mesure donnait-elle à peu près la même mémoire pour tous les moteurs~?
\end{enumerate}

\medskip\noindent\textbf{Partie IV --- Le verdict.}
\begin{enumerate}[resume]
\item Énoncer le \emph{résultat nommé}~: les écarts de volume et de rendement entre les conventions, ainsi que le temps et la mémoire du pipeline dans
les cinq modes.
\item Quel moteur retenir pour un mois, et lequel pour cinq ans des mêmes données~?
\item Que standardiser dans une équipe de recherche pour que les barres de deux chercheurs concordent~?
\item Contre quoi un test d'équivalence entre moteurs protège-t-il~?
\item En une phrase~: qu'est-ce qui rend un pipeline de \omterm{def:pl:dataframes:dataframe}{dataframes} fiable~?
\end{enumerate}
\end{problem}

\section{Questions d'entretien}

\begin{interviewq}[$\star$ \iqroles{chercheur}]\label{iq:pl:dataframes:1}
Vos barres de cinq minutes ne concordent pas avec celles d'un collègue. Que vérifiez-vous~?
\end{interviewq}

\begin{interviewq}[$\star\star$ \iqroles{chercheur, développeur}]\label{iq:pl:dataframes:2}
Quelle est la différence entre l'\omterm{def:pl:dataframes:eager}{évaluation immédiate} et l'\omterm{def:pl:dataframes:eager}{évaluation paresseuse} des \omterm{def:pl:dataframes:dataframe}{dataframes}, et dans quels cas cela importe-t-il~?
\end{interviewq}

\begin{interviewq}[$\star\star$ \iqroles{chercheur}]\label{iq:pl:dataframes:3}
Comment calculeriez-vous un rang en coupe transversale des rendements à une minute sur 3\,000 symboles pendant un an~?
\end{interviewq}

\begin{interviewq}[$\star\star$ \iqroles{développeur}]\label{iq:pl:dataframes:4}
Votre tâche pandas manque de mémoire sur un mois de transactions. Que modifiez-vous~?
\end{interviewq}

\begin{interviewq}[$\star\star\star$ \iqroles{développeur, chercheur}]\label{iq:pl:dataframes:5}
Concevez une bibliothèque partagée de transformations de données de marché pour une équipe de recherche utilisant plusieurs moteurs de \omterm{def:pl:dataframes:dataframe}{dataframes}.
\end{interviewq}
